% Shared preamble for the Week 7 student pages (Take-away games)
\documentclass[12pt]{article}
\usepackage[letterpaper, top=0.85in, bottom=0.8in, left=0.75in, right=0.75in,
            headheight=18pt, headsep=16pt, footskip=30pt]{geometry}
\usepackage{fontspec}
\setmainfont{Andika}
\usepackage{fancyhdr}
\usepackage{tikz}
\usetikzlibrary{shapes.geometric, calc}
\usepackage{array}

\pagestyle{fancy}
\fancyhf{}
\renewcommand{\headrulewidth}{0.5pt}
\renewcommand{\footrulewidth}{0pt}
\fancyfoot[L]{\fontsize{11}{13}\selectfont Bellingham Math Circle / Week 7}
\fancyfoot[R]{\fontsize{11}{13}\selectfont \thepage}
\newcommand{\header}[1]{\fancyhead[L]{\fontsize{12}{14}\selectfont\bfseries Week 7 / Take-away games / #1}}

\setlength{\parindent}{0pt}
\setlength{\parskip}{0pt}
\newcommand{\problem}[1]{\par\textbf{Problem #1:}~}
\hyphenpenalty=10000
\exhyphenpenalty=10000
\AtBeginDocument{\raggedright}

% ---------- drawing styles ----------
\tikzset{
  ctr/.style={circle, draw=black, line width=0.9pt, fill=black!22, inner sep=0pt, minimum size=#1},
  ctr/.default=0.6cm,
  wctr/.style={circle, draw=black, line width=0.9pt, fill=white, inner sep=0pt, minimum size=#1},
  wctr/.default=0.7cm,
  pbox/.style={draw=black!55, line width=0.8pt, rounded corners=7pt},
  estar/.style={star, star points=5, star point ratio=2.3, draw=black, line width=0.9pt,
                minimum size=#1, inner sep=0pt},
  estar/.default=1.15cm,
  choice/.style={draw=black, line width=0.8pt, rounded corners=8pt, minimum width=1.45cm,
                 minimum height=0.85cm, font=\fontsize{15}{15}\selectfont},
}

% counters in rows of five, first counter centred at (#2,#3), spacing #4, node style #5
\newcommand{\rowsoffive}[5]{%
  \foreach \ci in {1,...,#1}{%
    \pgfmathtruncatemacro{\prow}{(\ci-1)/5}%
    \pgfmathtruncatemacro{\pcol}{\ci-1-5*\prow}%
    \node[#5] at ({#2+\pcol*#4},{#3-\prow*#4}) {};}}

% a boxed pile of #1 counters centred at (#2,#3); box #4 x #5; spacing #6; style #7
\newcommand{\pileinbox}[7]{%
  \draw[pbox] ({#2-#4/2},{#3-#5/2}) rectangle ({#2+#4/2},{#3+#5/2});
  \pgfmathtruncatemacro{\pcols}{min(#1,5)}%
  \pgfmathtruncatemacro{\prows}{ceil(#1/5)}%
  \pgfmathsetmacro{\px}{#2-(\pcols-1)*#6/2}%
  \pgfmathsetmacro{\py}{#3+(\prows-1)*#6/2}%
  \rowsoffive{#1}{\px}{\py}{#6}{#7}}

% Standalone pile pictures: #1 n, #2 box width, #3 box height, #4 spacing, #5 counter style
\newcommand{\pilestar}[5]{%
  \begin{tikzpicture}[baseline=(current bounding box.north)]
    \pileinbox{#1}{0}{0}{#2}{#3}{#4}{#5}
    \node[font=\fontsize{15}{15}\selectfont] at (0,{-#3/2-0.42}) {#1};
    \node[estar] at (0,{-#3/2-1.35}) {};
  \end{tikzpicture}}
\newcommand{\pilechoice}[5]{%
  \begin{tikzpicture}[baseline=(current bounding box.north)]
    \pileinbox{#1}{0}{0}{#2}{#3}{#4}{#5}
    \node[font=\fontsize{15}{15}\selectfont] at (0,{-#3/2-0.42}) {#1};
    \node[choice] at (-0.95,{-#3/2-1.35}) {1st};
    \node[choice] at (0.95,{-#3/2-1.35}) {2nd};
  \end{tikzpicture}}
\newcommand{\pileplain}[5]{%
  \begin{tikzpicture}[baseline=(current bounding box.north)]
    \pileinbox{#1}{0}{0}{#2}{#3}{#4}{#5}
    \node[font=\fontsize{15}{15}\selectfont] at (0,{-#3/2-0.42}) {#1};
  \end{tikzpicture}}

% Two rows of counters, each row in its own tray, inside one box.
% #1 top row, #2 bottom row, #3 box width, #4 box height, #5 spacing, #6 distance between
% row centres, #7 counter style, #8 what goes underneath (code, y origin at box bottom)
\newcommand{\tworowpic}[8]{%
  \begin{tikzpicture}[baseline=(current bounding box.north)]
    \draw[pbox] ({-#3/2},{-#4/2}) rectangle ({#3/2},{#4/2});
    \pgfmathsetmacro{\trayw}{#3-0.5}%
    \pgfmathsetmacro{\trayh}{#5*0.95}%
    \pgfmathsetmacro{\lx}{-\trayw/2+0.25+#5/2}%
    \draw[black!45, line width=0.6pt, rounded corners=5pt]
      ({-\trayw/2},{#6/2-\trayh/2}) rectangle ({\trayw/2},{#6/2+\trayh/2});
    \draw[black!45, line width=0.6pt, rounded corners=5pt]
      ({-\trayw/2},{-#6/2-\trayh/2}) rectangle ({\trayw/2},{-#6/2+\trayh/2});
    \foreach \ci in {1,...,#1}{\node[#7] at ({\lx+(\ci-1)*#5},{#6/2}) {};}
    \foreach \ci in {1,...,#2}{\node[#7] at ({\lx+(\ci-1)*#5},{-#6/2}) {};}
    \begin{scope}[shift={(0,{-#4/2})}] #8 \end{scope}
  \end{tikzpicture}}

% Recording strip: numbers #1..#2, #3 cells per row, cell width #4, cell height #5
\newcommand{\recstrip}[5]{%
  \begin{tikzpicture}
    \foreach \k in {#1,...,#2}{%
      \pgfmathtruncatemacro{\srow}{(\k-#1)/#3}%
      \pgfmathtruncatemacro{\scol}{\k-#1-#3*\srow}%
      \pgfmathsetmacro{\sx}{\scol*#4}%
      \pgfmathsetmacro{\sy}{-\srow*(#5+0.25)}%
      \draw[line width=0.7pt] (\sx,\sy) rectangle ({\sx+#4},{\sy-#5});
      \node[font=\fontsize{13}{13}\selectfont\bfseries, anchor=base] at ({\sx+#4/2},{\sy-0.48}) {\k};
      \draw[black!35, line width=0.4pt] ({\sx+0.12},{\sy-0.66}) -- ({\sx+#4-0.12},{\sy-0.66});
    }
  \end{tikzpicture}}

% Strip of numbers to circle: #1..#2, #3 per row, spacing #4, row spacing #5
\newcommand{\circlestrip}[5]{%
  \begin{tikzpicture}
    \foreach \k in {#1,...,#2}{%
      \pgfmathtruncatemacro{\srow}{(\k-#1)/#3}%
      \pgfmathtruncatemacro{\scol}{\k-#1-#3*\srow}%
      \node[font=\fontsize{14}{14}\selectfont, anchor=base] at ({\scol*#4},{-\srow*#5}) {\k};
    }
  \end{tikzpicture}}

\header{Grades 2–3}

\newcommand{\bsize}{\fontsize{13.5}{19}\selectfont}
\newcommand{\twenty}{\begin{center}\recstrip{1}{20}{10}{1.72}{1.8}\end{center}}
\newcommand{\blank}[1]{\rule[-0.15em]{#1}{0.5pt}}
% two rows of counters with an answer line underneath
\newcommand{\trp}[2]{\tworowpic{#1}{#2}{7.0}{2.6}{0.92}{1.12}{ctr=0.66cm}{%
    \node[anchor=west, font=\fontsize{13}{13}\selectfont] at (-3.5,-0.75) {First or second?};
    \draw[line width=0.5pt] (0.4,-0.95) -- (3.5,-0.95);}}

\begin{document}
\bsize

In every game, two players take turns taking counters. Each problem says how many counters you may take on a turn. Whoever takes the last counter wins. A player can \emph{always win} if they can win no matter how the other player plays.

\vspace{9mm}
\problem{1} You may take 1 or 2 counters on each turn. For each starting pile from 1 to 20, find out who can always win: the player who goes first, or the player who goes second. Write 1st or 2nd under the number.

\vspace{4mm}
\twenty

% ------------------------------------------------------------------
\newpage
\problem{2} You may still take 1 or 2 counters on each turn. For each pile in the table, would you rather go first or second? If you would go first, how many counters would you take on your first turn?

\vspace{6mm}
\begin{center}
\renewcommand{\arraystretch}{1.0}
\begin{tabular}{|>{\centering\arraybackslash}m{2.6cm}|>{\centering\arraybackslash}m{5.0cm}|>{\centering\arraybackslash}m{6.4cm}|}
\hline
\rule[-0.35cm]{0pt}{1.0cm}\textbf{Pile} & \rule[-0.35cm]{0pt}{1.0cm}\textbf{First or second?} & \rule[-0.35cm]{0pt}{1.0cm}\textbf{How many I take first} \\
\hline
\rule[-0.6cm]{0pt}{1.5cm}24 counters & & \\ \hline
\rule[-0.6cm]{0pt}{1.5cm}29 counters & & \\ \hline
\rule[-0.6cm]{0pt}{1.5cm}31 counters & & \\ \hline
\rule[-0.6cm]{0pt}{1.5cm}45 counters & & \\ \hline
\end{tabular}
\end{center}

\vspace{6mm}
Then pick one of these piles. Explain why your way of playing wins no matter what your partner does.

% ------------------------------------------------------------------
\newpage
\problem{3} Now you may take 1, 3, or 4 counters on each turn. You may not take 2. For each starting pile from 1 to 20, who can always win? Write 1st or 2nd under the number.

\vspace{4mm}
\twenty

% ------------------------------------------------------------------
\newpage
\problem{4} The rule is still: take 1, 3, or 4 counters on each turn. Here is what five children say about this game. Decide whether each one is right, and explain how you know.

\newcommand{\claim}[3]{\vspace{5mm}\par\hangindent=2.1em\hangafter=0
  \makebox[0pt][r]{#1\hspace{0.6em}}#2: “#3”\par\vspace{30mm}}
\claim{(a)}{Mia}{With 5 counters, I go first and take 4. Then I can always win.}
\claim{(b)}{Leo}{With 7 counters, I go second, and I can always win.}
\claim{(c)}{Zoe}{With 9 counters, I go first, and I can always win.}
\claim{(d)}{Sam}{With 10 counters, I go first and take 1. Then I can always win.}
\claim{(e)}{Ana}{When you are allowed to take 4, taking 4 is always the best move.}

% ------------------------------------------------------------------
\newpage
\problem{5} The rule is still: take 1, 3, or 4 counters on each turn. Who can always win when the pile starts with 30 counters? With 50 counters? With 100 counters? Explain how you know.

\vspace{8mm}
\hspace*{2.1em}\makebox[3.6cm][l]{30 counters:}\blank{4cm}

\vspace{6mm}
\hspace*{2.1em}\makebox[3.6cm][l]{50 counters:}\blank{4cm}

\vspace{6mm}
\hspace*{2.1em}\makebox[3.6cm][l]{100 counters:}\blank{4cm}

% ------------------------------------------------------------------
\newpage
\problem{6} Now the counters are in two rows. On your turn, take 1 or 2 counters from one row (not from both rows). For each start below, would you rather go first or second?

\vspace{6mm}
\begin{center}
\trp{4}{4}\hfill\trp{3}{5}

\vspace{8mm}
\trp{3}{6}\hfill\trp{1}{4}
\end{center}

\vspace{8mm}
Then pick one start where you would go second. Explain how you win no matter what your partner does.

% ------------------------------------------------------------------
\newpage
\problem{7} For each rule, who can always win when the pile starts with 35 counters?

\vspace{6mm}
\newcommand{\rl}[2]{\par\hspace*{2.1em}\makebox[1.8em][l]{#1}\makebox[7.2cm][l]{#2}Who can always win? \blank{3.2cm}\par\vspace{7mm}}
\rl{(a)}{Take 1, 2, or 3 counters.}
\rl{(b)}{Take 1 or 3 counters.}
\rl{(c)}{Take 1, 2, 3, or 4 counters.}
\rl{(d)}{Take 1 or 4 counters.}

\end{document}
